\documentclass[12pt,a4paper]{article}

\usepackage{amsmath, amsthm, amssymb}
\usepackage{mathtools}
\usepackage{hyperref}
\usepackage{cleveref}
\usepackage{enumitem}

% ---------- Theorem environments ----------
\newtheorem{theorem}{Theorem}[section]
\newtheorem{proposition}[theorem]{Proposition}
\newtheorem{lemma}[theorem]{Lemma}
\newtheorem{corollary}[theorem]{Corollary}
\theoremstyle{definition}
\newtheorem{assumption}[theorem]{Assumption}
\newtheorem{remark}[theorem]{Remark}

% ---------- Notation ----------
\newcommand{\Ish}{I_\delta}
\newcommand{\NSh}{N_{\mathrm{Sh}}}
\newcommand{\KN}{K_N}
\newcommand{\KBL}{K_{\mathrm{BL}}}
\newcommand{\wcl}{\omega_n^{\mathrm{cl}}}
\newcommand{\rhocl}{\rho_n^{\mathrm{cl}}}
\newcommand{\rWKB}{r_n^{\mathrm{WKB}}}
\newcommand{\Eout}{\mathcal{E}_{\mathrm{out}}}

\title{%
  WKB Cover, Schur Control, and Unconditional Bulk Stability\\
  of Prolate Christoffel--Darboux Projectors%
}
\author{[Author]}
\date{May 2026}

\begin{document}
\maketitle

\begin{abstract}
We prove two complementary results for prolate spheroidal wave
functions (PSWFs) in the bulk regime $N \le \gamma\,N_{\mathrm{Sh}}$.

\textbf{(I) Schur--diagonal bound.}
The Christoffel--Darboux kernel $K_N(x,y) = \sum_{n<N}\psi_n(x)\psi_n(y)$
satisfies a uniform diagonal estimate
$K_N(x,x) \le C_{\delta,\gamma}\,Nc/(\pi T)$ on bulk sub-intervals
$I_\delta = [-(1-\delta)T,(1-\delta)T]$, via a kernel splitting
and a Schur-type argument using the classical band-limited kernel
$K_{\mathrm{BL}}$.

\textbf{(II) Bridge Lemma (WKB geometric cover).}
For $\gamma < 1/2$, the out-of-band spectral energy of PSWF
products satisfies
\[
  \Eout(f_{mn})
  := \int_{|\xi|>\omega} |\widehat{\psi_m\psi_n}(\xi)|^2\,d\xi
  \;\le\; C_{\gamma}\,e^{-\alpha_\gamma c}
\]
unconditionally for all $m,n \le \gamma\,N_{\mathrm{Sh}}$.
The proof rests on a purely geometric observation:
for $\gamma < 1/2$, the WKB forbidden zones of $\psi_m$ and
$\psi_n$ together cover the entire convolution overlap interval,
so at least one factor is exponentially suppressed at every
integration point.

Together, these results establish Assumption~\textup{ass:bulkconv}
of \cite{PaperIII} unconditionally for $\gamma < 1/2$,
completing that leg of the bulk stability program
initiated in \cite{PaperI,PaperII,PaperIII,PaperIV}.
The proof uses no microlocal machinery, no phase-cancellation
argument, and no individual eigenvalue estimate in the
Airy transition region $N \sim N_{\mathrm{Sh}}$.
\end{abstract}

\tableofcontents

%% =============================================================
\section{Introduction}
\label{sec:intro}
%% =============================================================

\subsection{Position in the bulk stability program}

This paper is the fifth in a series
\cite{PaperI,PaperII,PaperIII,PaperIV} establishing
quantitative spectral control of PSWF projectors in the bulk regime.
The central object is the Christoffel--Darboux kernel
\[
  K_N(x,y) := \sum_{n=0}^{N-1}\psi_n^{(c)}(x)\,\psi_n^{(c)}(y),
  \qquad N \le \gamma\,\NSh,\quad \NSh := \frac{2c}{\pi},
\]
and the associated spectral projector $P_N$ onto the PSWF
subspace $V_{N,c}$ of the Paley--Wiener space $\mathrm{PW}_\omega$
at bandwidth $\omega = c/T$.

In \cite{PaperIII}, the bulk tail bound
\[
  \|(I - P_N)f_{mn}\|_{L^2} \le C_\gamma\,e^{-\alpha_\gamma N},
  \qquad f_{mn} := \psi_m\psi_n,
\]
is reduced to two hypotheses:
\begin{itemize}
  \item[(H1)] \textit{Equidistribution:} $\psi_n^2 \approx \lambda_n\,\rhocl$
  uniformly on $I_\delta$ for $n \le \gamma\,\NSh$.
  Proved in \cite{PaperIV} via Pr\"ufer phase analysis.
  \item[(H2)] \textit{Assumption~ass:bulkconv:}
  $\Eout(f_{mn}) \le C_\gamma\,e^{-\alpha_\gamma c}$
  for all $m,n \le \gamma\,\NSh$.
  This was an open hypothesis in \cite{PaperIII}.
\end{itemize}
The present paper proves~(H2) unconditionally for $\gamma < 1/2$
via the Bridge Lemma (Section~\ref{sec:bridge}).

\subsection{The two main contributions}

\medskip\noindent
\textbf{Schur--diagonal control
(Sections~\ref{sec:diagonal}--\ref{sec:offdiag}).}
A kernel-splitting argument decomposes $K_N$ into the
explicit band-limited kernel $\KBL$ and exponentially small
error terms. This yields uniform diagonal and off-diagonal
estimates for $K_N$ on $I_\delta$, culminating in a Schur test
(Section~\ref{sec:schur}).

\medskip\noindent
\textbf{Bridge Lemma (Section~\ref{sec:bridge}).}
The key new result.
For $\gamma < 1/2$, the Fourier convolution structure of
$\widehat{f}_{mn} = \widehat\psi_m * \widehat\psi_n$ combined
with the WKB forbidden-zone $L^2$-concentration bound yields
$\Eout(f_{mn}) \le C_\gamma\,e^{-\alpha_\gamma c}$.
The proof is purely geometric: the two forbidden zones
$\{x > x_+(n)\}$ and $\{T - x > x_+(m)\}$ together
cover the full convolution overlap interval whenever $\gamma < 1/2$.
No oscillatory cancellation or phase-gap analysis is used.

\subsection{Why $\gamma < 1/2$ is the natural threshold}

The condition $\gamma < 1/2$ is sharp for the geometric cover:
\begin{itemize}
  \item If $\gamma < 1/2$: turning points satisfy
  $x_+(m) + x_+(n) \le 2\gamma T < T$, so the two
  forbidden zones cover $[0,T]$ completely.
  \item If $\gamma = 1/2$: the cover is exact with a
  degenerate boundary.
  \item If $\gamma > 1/2$: the interval
  $({\gamma T},\,(1-\gamma)T)$ is uncovered; both factors
  are oscillatory there, and a phase-cancellation argument
  (as in \cite{PaperI}) would be needed.
\end{itemize}
The extension to $\gamma \in [1/2, 1)$ remains open
(see Remark~\ref{rem:bridge:extension}).

\subsection{Separation of mechanisms}

This paper makes explicit a conceptual separation between
two independent decay mechanisms:
\begin{enumerate}
  \item \textbf{WKB tunneling} (this paper,
  Section~\ref{sec:bridge}): gives
  $\Eout(f_{mn}) \le Ce^{-\alpha c}$ uniformly in $m,n$;
  does not give decay in $|m-n|$.
  \item \textbf{Phase-gap / integration by parts}
  (\cite{PaperI}): gives algebraic off-diagonal decay
  $|E_{mn}| \le C(1+|m-n|)^{-p}$ for the discrete defect;
  does not use forbidden-zone geometry.
\end{enumerate}
The DSTP interface axiom of \cite{PaperI} encodes the
second mechanism and is entirely separate from~(H2).

\subsection{Organisation}

\begin{itemize}
  \item Section~\ref{sec:prelim}: notation and the Schur test.
  \item Section~\ref{sec:diagonal}: diagonal bound via
  Pr\"ufer amplitudes.
  \item Section~\ref{sec:offdiag}: off-diagonal bound
  via kernel splitting.
  \item Section~\ref{sec:bridge}: Bridge Lemma
  (WKB cover, exponential $\Eout$ decay).
  \item Section~\ref{sec:schur}: Schur test and proof
  of the main theorem.
\end{itemize}

%% =============================================================
\section{Preliminaries}
\label{sec:prelim}
%% =============================================================

\subsection{PSWF notation}

We use the Fourier convention
$\widehat{f}(\xi) = \int_{\mathbb{R}} f(x)\,e^{-2\pi i\xi x}\,dx$.
The bandwidth parameter is $c = \omega T > 0$,
where $\omega > 0$ is the half-bandwidth and
$T > 0$ is the half-time-interval.
The prolate spheroidal wave functions $\psi_n = \psi_n^{(c)}$
are the orthonormal eigenfunctions of the concentration
operator
$\mathcal{K}_c f(x)
= \int_{-T}^T \frac{\sin\omega(x-y)}{\pi(x-y)}f(y)\,dy$
on $L^2(\mathbb{R})$, with eigenvalues
$1 > \lambda_0(c) > \lambda_1(c) > \cdots > 0$
\cite{Slepian64,OsipovRokhlinXiao}.
Each $\psi_n$ is $\omega$-bandlimited and
$\|\psi_n\|_{L^2} = 1$.
We set $\NSh := 2c/\pi$ and work throughout in the bulk
regime $N \le \gamma\,\NSh$, $0 < \gamma < 1$.

\subsection{PSWF self-duality}
\label{subsec:selfduality}

\begin{equation}
\label{eq:selfduality}
  \widehat{\psi}_n(\xi)
  = (-i)^n\sqrt{\lambda_n}\,\psi_n\!\left(\frac{\xi T}{\omega}\right),
  \qquad |\xi| \le \omega
\end{equation}
(see \cite[eq.~(2.37)]{OsipovRokhlinXiao}),
and $\widehat\psi_n(\xi) = 0$ for $|\xi| > \omega$.

\subsection{Classical turning points}

For $n \le \gamma\,\NSh$, the classical turning point is
\begin{equation}
\label{eq:turning}
  x_+(n) := \frac{T\sqrt{\chi_n(c)}}{c} \;\le\; \gamma T + O(1/c)
\end{equation}
(see \cite[Theorem~2.7]{OsipovRokhlinXiao}).
The \emph{classical forbidden zone} of $\psi_n$ is
$(x_+(n), T)$.

\subsection{The Schur test}

\begin{lemma}[Schur test {\cite[Thm.~VI.23]{ReedSimonI}}]
\label{lem:schur}
Let $T_K$ be an integral operator with symmetric kernel
$K(x,y)$ on $L^2(I_\delta)$. If
\[
  S := \sup_{x \in I_\delta}\int_{I_\delta}|K(x,y)|\,dy < \infty,
\]
then $\|T_K\|_{L^2(I_\delta)\to L^2(I_\delta)} \le S$.
\end{lemma}

%% =============================================================
\section{The Diagonal Bound}
\label{sec:diagonal}
%% =============================================================

\begin{proposition}[Diagonal bound]
\label{prop:diagonal}
Fix $0 < \delta < 1$ and $0 < \gamma < 1$.
For all $x \in I_\delta$ and all $N \le \gamma\,\NSh$:
\begin{equation}
\label{eq:diagonal}
  K_N(x,x) = \sum_{n=0}^{N-1}\psi_n(x)^2
  \;\le\; C_{\delta,\gamma}\,\frac{Nc}{\pi T},
\end{equation}
uniformly in $c > 0$.
\end{proposition}

\begin{proof}
\textbf{Step~1 (Pr\"ufer amplitude).}
By \cite[Lemma~4.2]{PaperIV}:
\begin{equation}
\label{eq:amplitude-drift}
  r_n(x)^2 / \rWKB(x)^2 = 1 + O(1/n),
\end{equation}
and by \cite[Lemma~4.3]{PaperIV}:
\begin{equation}
\label{eq:normalization}
  \rWKB(x)^2/2 = \lambda_n(c)\,\rhocl(x)(1 + O(1/n)).
\end{equation}
Hence $\psi_n(x)^2 \le C_{\delta,\gamma}\,\lambda_n(c)\,\rhocl(x)$.

\textbf{Step~2 (Classical density).}
By \cite[Lemma~2.5(i)]{PaperIV}, $\wcl(x) \ge c_{\delta,\gamma}\,n$
on $I_\delta$, so the numerator satisfies $1/\wcl(x) \le C/n$.
The denominator satisfies $\int_{I_\delta} dt/\wcl(t)
\ge c'_{\delta,\gamma}T/c$. Therefore:
\begin{equation}
\label{eq:density-bound}
  \rhocl(x) \;\le\; C_{\delta,\gamma}\,\frac{c}{T}.
\end{equation}

\textbf{Step~3 (Summation).}
$K_N(x,x) \le C_{\delta,\gamma}(c/T)\sum_{n<N}\lambda_n
\le C_{\delta,\gamma}\,Nc/(\pi T)$.
\end{proof}

\begin{corollary}[$L^\infty$ bound]
\label{cor:Linfty}
For all $n \le \gamma\,\NSh$ and $x \in I_\delta$:
$|\psi_n(x)|^2 \le C_{\delta,\gamma}\,c/T$.
\end{corollary}

%% =============================================================
\section{Off-diagonal Decay via Kernel Splitting}
\label{sec:offdiag}
%% =============================================================

\subsection{The band-limited kernel}

\begin{equation}
\label{eq:KBL-def}
  \KBL(x,y) := \sum_{n=0}^{\infty}\lambda_n\,\psi_n(x)\psi_n(y)
  = \frac{\omega}{\pi}\cdot\frac{\sin\omega(x-y)}{\omega(x-y)}
\end{equation}
\cite[Prop.~3.1]{OsipovRokhlinXiao}.

\begin{lemma}[Band-limited kernel bound]
\label{lem:BL-bound}
For all $x,y \in \mathbb{R}$:
$|\KBL(x,y)| \le C/(1 + \omega|x-y|)$.
Moreover,
$\sup_{x \in I_\delta}\int_{I_\delta}|\KBL(x,y)|\,dy \le C_\delta$
uniformly in $c$.
\end{lemma}

\begin{proof}
Pointwise: $|\sin(u)/u| \le 2/(1+|u|)$.
Integral: substituting $v = \omega(x-y)$,
\[
  \int_{I_\delta}|\KBL(x,y)|\,dy
  = \frac{1}{\pi}\int_{-\omega \cdot 2T}^{\omega \cdot 2T}
    \left|\frac{\sin v}{v}\right|dv
  \;\le\;
  \frac{1}{\pi}\int_{-\infty}^{\infty}
    \left|\frac{\sin v}{v}\right|dv = 1,
\]
since $\int_{-\infty}^\infty |\mathrm{sinc}(v/\pi)|dv/\pi = 1$
is a classical result.
\end{proof}

\subsection{Eigenvalue concentration}

\begin{lemma}[Eigenvalue concentration]
\label{lem:eigenvalue-decay}
There exist $c_1, c_2 > 0$ depending only on $\gamma$ such that:
\begin{enumerate}[label=(\roman*)]
  \item For all $n \le \gamma\,\NSh$:
  $1 - \lambda_n(c) \le e^{-c_1 c}$.
  \item For all $n \ge \NSh$:
  $\lambda_n(c) \le e^{-c_2 c}$.
\end{enumerate}
\end{lemma}

\begin{proof}
(i) \cite[Theorem~4.6]{OsipovRokhlinXiao}: since
$c - n\pi/2 \ge c(1-\gamma) \ge c_\gamma c$ for
$n \le \gamma\NSh$.
(ii) \cite[Theorem~4.9]{OsipovRokhlinXiao}: since
$n\pi/2 - c \ge (\pi/2-1)c > 0$ for $n \ge \NSh$.
\end{proof}

\subsection{Kernel decomposition and error bounds}

Set $M := \lfloor\gamma\,\NSh\rfloor$. Write
\begin{equation}
\label{eq:split}
  K_M = \KBL + \widetilde{E}^- - \widetilde{E}^+,
\end{equation}
where $\widetilde{E}^- := \sum_{n<M}(1-\lambda_n)\psi_n\psi_n$
and $\widetilde{E}^+ := \sum_{n\ge M}\lambda_n\psi_n\psi_n$.
Split further $\widetilde{E}^+ = \widetilde{E}^+_{\mathrm{near}}
+ \widetilde{E}^+_{\mathrm{tail}}$ at $n = \NSh$.

\medskip\noindent
\textbf{Bound on $\widetilde{E}^-$.}
By Lemma~\ref{lem:eigenvalue-decay}(i) and
Corollary~\ref{cor:Linfty}:
\begin{equation}
\label{eq:Eminus-final}
  |\widetilde{E}^-(x,y)|
  \le C_{\delta,\gamma}\,\frac{c^2}{T}\,e^{-c_1 c}.
\end{equation}

\medskip\noindent
\textbf{Bound on $\widetilde{E}^+_{\mathrm{tail}}$.}
By Lemma~\ref{lem:eigenvalue-decay}(ii) and
Corollary~\ref{cor:Linfty}:
\begin{equation}
\label{eq:Eplus-tail-final}
  |\widetilde{E}^+_{\mathrm{tail}}(x,y)|
  \le C_{\delta,\gamma}\,\frac{c}{T}\,e^{-c_2 c}.
\end{equation}

\medskip\noindent
\textbf{Bound on $\widetilde{E}^+_{\mathrm{near}}$
via PSD dominance.}
Since $0 \preceq \widetilde{E}^+_{\mathrm{near}}
\preceq \KBL$, the Cauchy--Schwarz inequality for PSD
kernels gives $|A(x,y)|^2 \le A(x,x)A(y,y)$, hence:
\begin{equation}
\label{eq:Eplus-near-pointwise}
  |\widetilde{E}^+_{\mathrm{near}}(x,y)|
  \;\le\; \KBL(x,x)^{1/2}\KBL(y,y)^{1/2}
  = \frac{\omega}{\pi} = \frac{c}{\pi T}.
\end{equation}
No eigenvalue estimate in the Airy transition region
$M \le n < \NSh$ is required.

\begin{proposition}[Off-diagonal bound]
\label{prop:offdiag}
For all $x,y \in I_\delta$ and $N \le \gamma\,\NSh$:
\begin{equation}
\label{eq:offdiag-final}
  |K_N(x,y)|
  \;\le\;
  \frac{C_{\delta,\gamma}}{1+\omega|x-y|}
  + C_{\delta,\gamma}\,\frac{c}{\pi T}
  + C_{\delta,\gamma}\,\frac{c^2}{T}\,e^{-c_3 c},
\end{equation}
where $c_3 := \min(c_1,c_2) > 0$, uniformly in $c,T,N$.
\end{proposition}

\begin{proof}
Triangle inequality on~\eqref{eq:split}, then
Lemma~\ref{lem:BL-bound} and
\eqref{eq:Eminus-final}--\eqref{eq:Eplus-near-pointwise}.
\end{proof}

%% =============================================================
\section{The Bridge Lemma}
\label{sec:bridge}
%% =============================================================

\subsection{Statement}

For $m,n \le \gamma\,\NSh$, set $f_{mn} := \psi_m\psi_n$ and
\[
  \Eout(f_{mn}) := \int_{|\xi|>\omega}|\widehat{f}_{mn}(\xi)|^2\,d\xi.
\]

\begin{theorem}[Bridge Lemma]
\label{thm:bridge}
Fix $0 < \gamma < 1/2$. There exist $C_\gamma, \alpha_\gamma > 0$
such that for all $m,n \le \gamma\,\NSh$ and $c \ge 1$:
\begin{equation}
\label{eq:bridge:main}
  \Eout(f_{mn}) \;\le\; C_\gamma\,e^{-\alpha_\gamma c}.
\end{equation}
This establishes Assumption~ass:bulkconv of \cite{PaperIII}
unconditionally for $\gamma < 1/2$.
\end{theorem}

\subsection{Key input: $L^2$-concentration}

\begin{lemma}[$L^2$-forbidden-zone concentration]
\label{lem:L2-forbidden}
For all $n \le \gamma\,\NSh$ and $0 \le a < T$:
\begin{equation}
\label{eq:L2-forbidden}
  \|\psi_n\|_{L^2[a,T]}^2
  \;\le\; 1 - \lambda_n(c) \;\le\; e^{-c_1 c}.
\end{equation}
\end{lemma}

\begin{proof}
$\|\psi_n\|_{L^2[a,T]}^2 \le \|\psi_n\|_{L^2(\mathbb{R}\setminus[-T,T])}^2
= 1 - \lambda_n$. Apply Lemma~\ref{lem:eigenvalue-decay}(i).
\end{proof}

\subsection{The geometric cover}

\begin{lemma}[Geometric forbidden-zone cover]
\label{lem:bridge:cover}
Fix $\gamma < 1/2$, $m,n \le \gamma\,\NSh$, and
$x_0 = T + \delta T/\omega$ for $\delta \in (0,\omega)$.
For every $x \in [\delta T/\omega, T]$, at least one holds:
\begin{enumerate}[label=(\Roman*)]
  \item $x \ge \gamma T$: $\psi_n(x)$ is in its forbidden zone.
  \item $x < \gamma T$: $x_0 - x \ge (1-\gamma)T > \gamma T
  \ge x_+(m)$, so $\psi_m(x_0-x)$ is in its forbidden zone.
\end{enumerate}
\end{lemma}

\begin{proof}
If $x < \gamma T$: $x_0 - x > T - \gamma T = (1-\gamma)T$.
Since $\gamma < 1/2$: $(1-\gamma)T > \gamma T \ge x_+(m)$.
\end{proof}

\begin{remark}[Sharpness]
\label{rem:bridge:halfsharp}
For $\gamma \ge 1/2$, the interval
$x \in (\gamma T, (1-\gamma)T)$ is non-empty and neither
case applies there: both factors are oscillatory and the
geometric cover fails.
\end{remark}

\subsection{Proof of Theorem~\ref{thm:bridge}}

\begin{proof}
\textbf{Step~1 (Convolution representation).}
By~\eqref{eq:selfduality} and the substitution $x = \eta T/\omega$,
for $\xi = \omega+\delta$, $\delta \in (0,\omega)$:
\begin{equation}
\label{eq:bridge:conv}
  \widehat{f}_{mn}(\omega+\delta)
  = (-i)^{m+n}\sqrt{\lambda_m\lambda_n}\,\frac{\omega}{T}
    \int_{\delta T/\omega}^{T}
    \psi_m\!\bigl(T + \tfrac{\delta T}{\omega} - x\bigr)\,
    \psi_n(x)\,dx.
\end{equation}

\textbf{Step~2 (Splitting).}
Set $x_0 = T + \delta T/\omega$ and split at $\gamma T$:
\[
  \int_{\delta T/\omega}^{T}\psi_m(x_0-x)\psi_n(x)\,dx
  = I_{\mathrm{I}} + I_{\mathrm{II}},
\]
where $I_{\mathrm{I}}$ integrates over $[\gamma T, T]$
and $I_{\mathrm{II}}$ over $[\delta T/\omega, \gamma T]$.

\textbf{Step~3 (Bound on $I_{\mathrm{I}}$).}
For $x \in [\gamma T, T]$: $\psi_n(x)$ is in its
forbidden zone. By Cauchy--Schwarz and
Lemma~\ref{lem:L2-forbidden}:
\begin{equation}
\label{eq:bridge:I1}
  |I_{\mathrm{I}}|
  \le \|\psi_m\|_{L^\infty}\,\|\psi_n\|_{L^1[\gamma T,T]}
  \le C\,T^{1/2}\,\|\psi_n\|_{L^2[\gamma T,T]}
  \le C\,T^{1/2}\,e^{-c_1 c/2},
\end{equation}
using $\|\psi_m\|_{L^\infty} \le C$
\cite[Corollary~4.3]{OsipovRokhlinXiao}.

\textbf{Step~4 (Bound on $I_{\mathrm{II}}$).}
For $x \in [\delta T/\omega, \gamma T]$:
by Lemma~\ref{lem:bridge:cover}, $x_0 - x \ge (1-\gamma)T$
lies in the forbidden zone of $\psi_m$.
Substituting $u = x_0 - x$:
\begin{equation}
\label{eq:bridge:I2}
  |I_{\mathrm{II}}|
  \le \|\psi_n\|_{L^\infty}\,\|\psi_m\|_{L^1[(1-\gamma)T,T]}
  \le C\,T^{1/2}\,\|\psi_m\|_{L^2[(1-\gamma)T,T]}
  \le C\,T^{1/2}\,e^{-c_1 c/2}.
\end{equation}

\textbf{Step~5 (Pointwise Fourier bound).}
From~\eqref{eq:bridge:conv} and Steps~3--4,
using $\sqrt{\lambda_m\lambda_n} \le 1$:
\begin{equation}
\label{eq:bridge:pointwise}
  |\widehat{f}_{mn}(\omega+\delta)|
  \;\le\; \frac{2C\omega}{T^{1/2}}\,e^{-c_1 c/2}.
\end{equation}

\textbf{Step~6 ($\Eout$ bound).}
\begin{equation}
\label{eq:bridge:Eout}
  \Eout(f_{mn})
  = 2\int_0^\omega |\widehat{f}_{mn}(\omega+\delta)|^2\,d\delta
  \;\le\; \frac{8C^2\omega^3}{T}\,e^{-c_1 c}.
\end{equation}
Since $\omega^3/T = c^3/T^4$ and $c^3 e^{-c_1 c}
\le C_\varepsilon\,e^{-(c_1-\varepsilon)c}$ for any
$\varepsilon \in (0,c_1)$, setting
$\alpha_\gamma := c_1 - \varepsilon$ yields~\eqref{eq:bridge:main}.
\end{proof}

\subsection{Consequences}

\begin{corollary}[ass:bulkconv for $\gamma < 1/2$]
\label{cor:bulkconv}
For all $0 < \gamma < 1/2$ and $m,n \le \gamma\,\NSh$,
Assumption~ass:bulkconv of \cite{PaperIII} holds.
By \cite[Theorem~3.2]{PaperIII}:
\[
  \|(I-P_N)f_{mn}\|_{L^2}
  \;\le\; C_\gamma\,e^{-\alpha_\gamma' N}
\]
holds unconditionally for all $m,n,N \le \gamma\,\NSh$.
\end{corollary}

\begin{remark}[Extension to $\gamma \ge 1/2$]
\label{rem:bridge:extension}
For $\gamma \in [1/2,1)$, the interval
$x \in (\gamma T,(1-\gamma)T)$ is uncovered.
Both factors are oscillatory there; an
integration-by-parts argument on the WKB phase
$\Phi(x) = \phi_m(x_0-x) \pm \phi_n(x)$ would be needed,
analogous to \cite{PaperI}.
Whether $\Eout(f_{mn}) \le Ce^{-\alpha c}$ extends to
$\gamma \ge 1/2$ remains open.
\end{remark}

\begin{remark}[Separation of mechanisms]
\label{rem:bridge:separation}
The Bridge Lemma uses only:
(i) $L^2$ eigenvalue concentration
(Lemma~\ref{lem:L2-forbidden}), and
(ii) the geometric cover for $\gamma < 1/2$.
No phase-gap estimate, no integration by parts, and
no index-dependent decay in $|m-n|$ is used.
The DSTP off-diagonal decay
$|E_{mn}| \le C(1+|m-n|)^{-p}$ of \cite{PaperI}
requires the separate phase-gap mechanism of \cite{PaperI}.
The two mechanisms are logically independent.
\end{remark}

%% =============================================================
\section{The Schur Argument}
\label{sec:schur}
%% =============================================================

\begin{theorem}[Diagonal Schur bound]
\label{thm:schur-main}
Fix $0 < \delta < 1$ and $0 < \gamma < 1$.
For all $N \le \gamma\,\NSh$ and $c \ge 1$:
\begin{equation}
\label{eq:schur-main}
  \sup_{x \in I_\delta} K_N(x,x)
  \;\le\; C_{\delta,\gamma}\,\frac{Nc}{\pi T}.
\end{equation}
Moreover, the Schur integrability condition holds:
\begin{equation}
\label{eq:schur-cond}
  \sup_{x \in I_\delta}
  \int_{I_\delta} |K_N(x,y)|\,dy
  \;\le\; C_{\delta,\gamma}
  \left(1 + \frac{2(1-\delta)c}{\pi}\right).
\end{equation}
\end{theorem}

\begin{remark}[Interpretation]
\label{rem:schur-interp}
The bound~\eqref{eq:schur-cond} grows like $c/\pi$,
reflecting the spectral density $\sim Nc/T$ of $K_N$.
This is expected and does not contradict the Bridge Lemma:
the uniform bound on $\Eout$ (Theorem~\ref{thm:bridge})
is a Fourier-side statement and provides the exponential
correction that makes the bulk tail bound
of \cite[Theorem~3.2]{PaperIII} unconditional.
The Schur bound here provides the polynomial diagonal
control used in that argument.
\end{remark}

\begin{proof}[Proof of Theorem~\ref{thm:schur-main}]
\textbf{Bound~\eqref{eq:schur-main}}: Proposition~\ref{prop:diagonal}.

\textbf{Bound~\eqref{eq:schur-cond}:}
Integrate Proposition~\ref{prop:offdiag} over $y \in I_\delta$.

\textit{Sinc term:} Substituting $v = \omega(x-y)$,
\[
  \int_{I_\delta}\frac{C_{\delta,\gamma}}{1+\omega|x-y|}\,dy
  = \frac{C_{\delta,\gamma}}{\omega}
    \int_{-2\omega T}^{2\omega T}\frac{dv}{1+|v|}
  \le \frac{2C_{\delta,\gamma}}{\omega}
    \log(1+2\omega T).
\]
Since $\omega = c/T$: $(2/\omega)\log(1+2\omega T)
= (2T/c)\log(1+2c/T) \le C'_{\delta,\gamma}$
for all $c \ge 1$, $T$ fixed.

\textit{Near-bulk term:}
$\int_{I_\delta}(c/\pi T)\,dy = 2(1-\delta)c/\pi$.

\textit{Exponential tail:}
$\int_{I_\delta}(c^2/T)e^{-c_3 c}\,dy
\le 2c^2 e^{-c_3 c} \le C'_{\delta,\gamma}$.

Summing gives~\eqref{eq:schur-cond}.
\end{proof}

\begin{corollary}[Unconditional bulk tail bound]
\label{cor:bulk-final}
For all $0 < \gamma < 1/2$ and $m,n,N \le \gamma\,\NSh$:
\[
  \|(I - P_N)f_{mn}\|_{L^2([-T,T])}
  \;\le\; C_\gamma\,e^{-\alpha_\gamma' N}.
\]
\end{corollary}

\begin{proof}
Combine Corollary~\ref{cor:bulkconv} with
\cite[Theorem~3.2]{PaperIII}: both hypotheses (H1) and (H2)
are now established, (H1) by \cite{PaperIV} and (H2)
by Theorem~\ref{thm:bridge}.
\end{proof}

%% =============================================================
\begin{thebibliography}{99}

\bibitem{PaperI}
[Author], \textit{Gram Coercivity of PSWF Quadrature Forms},
preprint, 2025.

\bibitem{PaperII}
[Author], \textit{Scaling Limits, Trace Formula, and Weil
Connection}, preprint, 2025.

\bibitem{PaperIII}
[Author], \textit{PSWF Product Spectral Tail Estimates},
preprint, 2026.

\bibitem{PaperIV}
[Author], \textit{Semiclassical Equidistribution of PSWF
Densities}, preprint, 2026.

\bibitem{OsipovRokhlinXiao}
A.~Osipov, V.~Rokhlin, and H.~Xiao,
\textit{Prolate Spheroidal Wave Functions of Order Zero},
Springer, 2013.

\bibitem{ReedSimonI}
M.~Reed and B.~Simon,
\textit{Methods of Modern Mathematical Physics, Vol.~I},
Academic Press, 1972.

\bibitem{Slepian64}
D.~Slepian,
\textit{Prolate spheroidal wave functions, Fourier analysis
and uncertainty~-- IV},
Bell Syst.\ Tech.\ J.\ \textbf{43} (1964), 3009--3057.

\bibitem{Widom64}
H.~Widom,
\textit{Asymptotic behavior of the eigenvalues of certain
integral equations~II},
Arch.\ Rational Mech.\ Anal.\ \textbf{17} (1964), 215--229.

\end{thebibliography}

\end{document}
